\section{An exact telescoping antiderivative calculus}\label{sec:primitives}

\begin{theorem}[Reference-parameter telescoping]\label{thm:primitive}
For $K\ge1$ and $\Rea u>-K$,
\begin{equation}\label{eq:primitiveDerivative}
 \boxed{\partial_c R_K(u;c)
 =(u+K)A_{K-1}(u;c)\zeta(1+u+K,c).}
\end{equation}
Consequently an exact primitive is
\begin{equation}\label{eq:primitive}
 \int A_{K-1}(u;c)\zeta(1+u+K,c)\,\dd c
 =-\frac1{u+K}\sum_{j=0}^{K-1}
 A_j(u;c)\zeta(1+u+j,c)+C.
\end{equation}
At a removable resonance, the finite sum on the right is evaluated as a
combined Laurent expression; any residual pole is independent of $c$ and
can be absorbed into the integration constant before taking the limit.
Equivalently, $R_K(u;c)/(u+K)$ is a primitive that is holomorphic throughout
the stated half-plane.
\end{theorem}
\begin{proof}
Differentiate the finite counterterm sum in \eqref{eq:master}, using
\eqref{eq:cderiv} and
$\partial_c\zeta(s,c)=-s\zeta(s+1,c)$.
For $j\ge1$ the derivative of the coefficient contributes
$(u+j)A_{j-1}\zeta(1+u+j,c)$, which cancels the derivative of the preceding
zeta factor.  Only
$-(u+K)A_{K-1}\zeta(1+u+K,c)$ remains from differentiating the sum.
The $c$-independent $G(u)$ contributes zero, proving
\eqref{eq:primitiveDerivative}.  Division by $u+K$ is harmless on
$\Rea u>-K$ and gives the primitive.  The last statement explains how to
read its meromorphic finite-sum form at a resonance.
\end{proof}

For $K=N+1$ and $u=-N+t$, the all-jet primitive law is particularly simple:
\begin{equation}\label{eq:primitivejets}
 \partial_c S_{N,m}^{(N+1)}(c)
 =[t^m]\{(1+t)A_N(-N+t;c)\zeta(2+t,c)\}.
\end{equation}
Thus the finite Stieltjes expression in \eqref{eq:alljets} is an explicitly
evaluated antiderivative of a polynomial-weighted combination of spectral
zeta derivatives at $2$.  This is stronger than merely leaving the
antiderivative as an unevaluated integral.

The first level recovers
\begin{equation}\label{eq:stieltjesderivative}
 \gamma_1'(c)=\zeta(2,c)+\zeta^{[1]}(2,c),
\end{equation}
by differentiating \eqref{eq:stieltjes} in $c$.
For $N=1$, let $\rho=-(1-a)(1-b)$ and
$e_2=[t^2]E_1(t)$.  Since
\[
 A_1(-1+t;c)=\rho+(c-a-b+\tfrac32)t-\tfrac12t^2,
\]
Theorem~\ref{thm:jets} becomes
\begin{align}
 S_{1,1}^{(2)}(c)={}&e_2+\frac12+\frac12\log(2\pi)
 -\log\Gamma(c)+(c-a-b+\tfrac32)\psi(c)+\rho\gamma_1(c).
 \label{eq:firstprimitiveclosed}
\end{align}
Taking the $c$ derivative yields the concrete integral identity
\begin{align}
 &\int\bigl\{(c-a-b+\tfrac32+\rho)\zeta(2,c)
                  +\rho\zeta^{[1]}(2,c)\bigr\}\,\dd c\notag\\
 &\hspace{8mm}=-\log\Gamma(c)+(c-a-b+\tfrac32)\psi(c)
                   +\rho\gamma_1(c)+C.
 \label{eq:explicitprimitive}
\end{align}
For $a=b=\tfrac12$ this reads
\begin{equation}\label{eq:halfprimitive}
 \int\left\{(c+\tfrac14)\psi'(c)-\tfrac14\zeta^{[1]}(2,c)\right\}\dd c
 =-\log\Gamma(c)+(c+\tfrac12)\psi(c)-\tfrac14\gamma_1(c)+C.
\end{equation}
The all-order theorem gives the analogous formulas with higher spectral
jets, rather than requiring separate integration-by-parts calculations.

At zeroth order, \eqref{eq:primitivejets} says
$\partial_c R_{N+1}(-N;c)=\rho_N\zeta(2,c)$.
Integrating from $c$ to $d$ gives the exact reference-change rule
\begin{equation}\label{eq:resreference}
 R_{N+1}(-N;d)-R_{N+1}(-N;c)=\rho_N\{\psi(d)-\psi(c)\}.
\end{equation}
This makes the subtraction convention fully visible.  The finite constant
is not invariant under replacing $(n+c)^{-1}$ by an unspecified harmonic
reference term.
